\section{Introducción}

En muchas ciudades, moverse en carro se ha vuelto cada vez más difícil. Cada año hay más vehículos, pero las vías y la infraestructura siguen casi iguales, y al mismo tiempo se intenta que el transporte genere menos contaminación y menos impacto en la ciudad. Una parte importante de esa congestión no se da solo mientras se conduce, sino también cuando las personas buscan dónde parquear y se encuentran con una gestión poco eficiente de los parqueaderos. En muchos casos todavía se utilizan procesos manuales o sistemas que no están conectados entre sí. Esto termina generando filas en las entradas, respuestas tardías sobre si hay cupos disponibles y muy poco uso de los datos históricos para entender cómo se están usando realmente estos espacios.

En este contexto, los sistemas de Reconocimiento Automático de Matrículas (\textit{Automatic Number Plate Recognition}, ANPR) se usan cada vez más para automatizar la entrada y salida de vehículos. En pocas palabras, una cámara lee la placa y el sistema registra el ingreso o la salida sin que la persona tenga que hacer un trámite manual. Cuando estos sistemas se integran con un parqueadero “inteligente”, es posible asociar las placas con usuarios o reservas y consultar con más detalle cómo se están usando los cupos a lo largo del tiempo.
Aun así, muchas de las soluciones ANPR que se ofrecen en el mercado siguen siendo de código cerrado, requieren licencias costosas y dejan al cliente muy dependiente de un proveedor específico. Además, no suelen adaptarse bien cuando se quieren usar en universidades, en proyectos de investigación o en prototipos que necesitan cambiar y mejorar con rapidez.

Al mismo tiempo, muchas organizaciones han empezado a trabajar con prácticas de \textit{DevOps} y con esquemas de integración y entrega continua (CI/CD), que han cambiado la forma de construir y publicar aplicaciones de software, especialmente en arquitecturas basadas en microservicios. La idea principal es automatizar lo más posible las pruebas, la compilación y el despliegue, de manera que los cambios en el código lleguen más rápido y con menos errores a producción~\cite{devops_tools_2019,ci_cd_quality_2025}. Sin embargo, cuando el sistema tiene que trabajar con imágenes y con modelos de inteligencia artificial más complejos, aplicar este tipo de automatización ya no es tan sencillo. En la práctica, estos modelos suelen ser más pesados, tardan más en ejecutarse y necesitan computadores más potentes (por ejemplo, con tarjeta gráfica dedicada), lo que hace mucho más difícil diseñar y mantener un proceso automático de pruebas y despliegue que sea estable.

En este trabajo se propone una arquitectura completa para un sistema ANPR enfocado en la gestión de parqueaderos, basada en tecnologías de código abierto y apoyada en un flujo de trabajo DevOps que facilita su despliegue y su evolución en el tiempo. El sistema se toma como caso de estudio bajo el nombre \textbf{ANPR-VISION}, pero la descripción se plantea de forma lo bastante general como para que pueda servir de guía en otros proyectos de parqueadero inteligente o de control de acceso vehicular.

\subsection{Problema de investigación}

A partir de lo que se ha encontrado en trabajos previos y de lo que se ve en el funcionamiento real de muchos parqueaderos, surge la siguiente pregunta central:

\begin{quote}
¿Cómo diseñar una arquitectura para un sistema ANPR, basada en tecnologías abiertas, que permita automatizar la gestión de parqueaderos y que, al mismo tiempo, pueda desplegarse y mantenerse de forma ágil usando prácticas de DevOps?
\end{quote}

Esta pregunta se conecta, en la práctica, con dos tipos de retos. El primero tiene que ver con la parte técnica de detectar y leer matrículas en condiciones reales, donde la luz cambia, la cámara no siempre está bien ubicada o los vehículos se mueven rápido. El segundo está relacionado con la operación del sistema: debe estar disponible todo el tiempo y poder crecer cuando aumenta la demanda, procesando video y eventos en tiempo real sin que el servicio se detenga.

\subsection{Objetivo general}

Diseñar y explicar una arquitectura de referencia para un sistema de Reconocimiento Automático de Matrículas orientado a la gestión inteligente de parqueaderos, usando herramientas de código abierto y prácticas de DevOps que hagan más sencillo su despliegue, actualización y mantenimiento.

\subsection{Objetivos específicos}

\begin{itemize}
    \item Reunir y comentar los trabajos más relevantes sobre reconocimiento automático de matrículas, parqueaderos inteligentes y control de acceso de vehículos.
    \item Explicar de manera clara qué partes debería tener un sistema ANPR para parqueaderos y cómo se comunican entre sí.
    \item Plantear una forma ordenada de construir, probar y poner en funcionamiento el sistema, apoyándose en herramientas de automatización.
    \item Desarrollar un prototipo que funcione y usarlo como ejemplo para mostrar qué tan viable es la propuesta.
    \item Señalar qué aporta la propuesta frente a otras soluciones, qué limitaciones tiene y qué se podría mejorar en trabajos futuros.
\end{itemize}

\subsection{Estructura del artículo}

El artículo se divide en varias secciones. En la Sección~\ref{sec:relacionados} se presentan los trabajos previos sobre ANPR, parqueaderos inteligentes y el uso de DevOps en este tipo de sistemas. La Sección~\ref{sec:metodologia} explica el enfoque que se siguió para plantear y concretar la arquitectura. En la Sección~\ref{sec:implementacion} se describe con más detalle la solución propuesta y el caso de estudio construido. La Sección~\ref{sec:resultados} reúne los principales resultados y los criterios usados para evaluarlos. Finalmente, la Sección~\ref{sec:conclusiones} recoge las conclusiones más importantes y sugiere algunas líneas posibles para trabajos futuros.